\documentclass[11pt,a4paper]{article}

\usepackage{main}
\usepackage{exercices}

\renewcommand{\typedocument}{TP}
\renewcommand{\classe}{5e}
\renewcommand{\titre}{Electricité - Compte-rendu de TP}

\begin{document}

\creertitre{Compte-rendu de TP 1: Circuits électriques}


\textbf{Consignes de sécurité}

\emph{La manipulation de matériel électrique nécessite de respecter certaines consignes afin d'éviter tout accident.\\[-0.5em]  
\begin{itemize}
    \item Vérifier les branchements avant d'allumer le générateur, ne pas manipuler tant qu'il est allumé.
    \item Ne pas laisser le générateur allumé inutilement.
    \item Si un composant chauffe, éteindre le générateur et prévenir le professeur.
\end{itemize}
}

\vspace{1cm}

\textbf{I. Matériel}
    
\begin{itemize}
    \item Un générateur
    \item Une lampe
    \item Un moteur
    \item Un interrupteur
    \item Des fils de connexion
    \item Des pinces crocodiles
\end{itemize}

\vspace{1cm}

\textbf{II. Réalisation d'un circuit électrique simple}

\emph{Dessiner les schémas à chaque nouveau montage.}

\begin{enumerate}[label=\alph*)]
    \item Brancher le générateur et le régler pour obtenir une tension de 3V.
    \item Le relier avec un fil de connexion au moteur. Qu'est-ce qui se passe ? Le moteur tourne-t-il ?
    \item Relier un deuxième fil de connexion du générateur au moteur. Que se passe-t-il maintenant ?
    \item Ajouter un interrupteur entre le générateur et le moteur. Que se passe-t-il lorsque l'interrupteur est ouvert ou fermé ?
    \item Quelles sont les deux conditions nécessaires pour que le moteur puisse fonctionner ?
\end{enumerate}

\vspace{1cm}

\textbf{III. Sens du courant}

\begin{enumerate}[label=\alph*)]
    \setcounter{enumi}{5}
    \item Noter dans quel sens les hélices du moteur tournent. Inverser les connexions du générateur et observer à nouveau le sens de rotation des hélices. Que constatez-vous ?
\end{enumerate}

\vspace{1cm}

\textbf{IV. Matériaux conducteurs et isolants}

\begin{enumerate}[label=\alph*)]
    \setcounter{enumi}{6}
    \item Remplacer l’interrupteur par les pinces crocodiles. Tester différents objets et remplir le tableau (règle, ciseaux, papier, mine de critérium, etc\dots).
\end{enumerate}

\renewcommand{\arraystretch}{2.5}
\begin{tabular}{|>{\bfseries\columncolor{headergrey}\centering\arraybackslash}m{5cm}|*{4}{>{\centering\arraybackslash}m{3cm}|}}
    \hline
    Matériau & Règle (plastique) & & & \\
    \hline
    Moteur allumé ? & Non & & & \\
    \hline
    Conducteur ou isolant ? & Isolant & & & \\
    \hline
\end{tabular}

\end{document}